\originalchapter{偶极辐射与运动电荷}\label{ch:radiation}
\originalsection{局域源的辐射区}
设源尺度 $a$、频率 $\omega$, 满足 $a\ll c/\omega$, 观察距离 $r\gg a$. 若只保留辐射项, 还要求 $r\gg c/\omega$, 即 $kr\gg1$, $k=\omega/c$. $a\ll r\ll c/\omega$ 虽远离源的几何尺寸, 仍属准静态近场. 推迟势核可取 $1/R\simeq1/r$、$t-R/c\simeq t-r/c+\nn\cdot\rr'/c$. 最低阶电偶极辐射还忽略源内部的相位差. 连续性方程给
\[
\dot\pp=\int\rr'\partial_t\rho\dd^3x'
=-\int\rr'\nabla'\cdot\jj\dd^3x'=\int\jj\dd^3x',
\]
表面项由源局域性消失, 故 $\vct A\simeq\mu_0\dot\pp(t-r/c)/(4\pi r)$. 势的微分中, 对 $1/r$ 求导给近场, 对推迟时间求导给 $1/r$ 辐射场. 再用辐射区的横向关系,
\[
\ee_\mathrm{rad}=\frac{\mu_0}{4\pi r}\nn\times[\nn\times\ddot\pp(t-r/c)],\qquad
\bb_\mathrm{rad}=\frac1c\nn\times\ee_\mathrm{rad}.
\]
横向投影可从电偶极势直接看出:
\[
V={1\over4\pi\eps}\left[\frac{\nn\cdot\pp}{r^2}+\frac{\nn\cdot\dot\pp}{cr}\right]_{t-r/c},\qquad
\vct A={\mu_0\over4\pi r}\dot\pp(t-r/c).
\]
只取 $1/r$ 项时, $-\nabla V$ 对推迟时间求导给 $\nn(\nn\cdot\ddot\pp)/(4\pi\eps c^2r)$, $-\partial_t\vct A$ 给 $-\ddot\pp/(4\pi\eps c^2r)$. 二者合为 $\nn\times(\nn\times\ddot\pp)$, 纵向分量相消. 对理想点偶极, 或各主要频率均有 $ka\ll1$ 的局域源, 该式在辐射区给电偶极主项; 不表示任意快速变化的有限源都可忽略内部相位.

辐射场的能流为 $S=E_\mathrm{rad}^2/(\mu_0c)$, 单位立体角功率为 $\dd P/\dd\Omega=r^2S$. 于是瞬时角功率为 $\mu_0|\ddot\pp|^2\sin^2\theta/(16\pi^2c)$. 用
$\int\sin^2\theta\dd\Omega=2\pi\int_0^\pi\sin^3\theta\dd\theta=8\pi/3$, 得 $P=|\ddot\pp|^2/(6\pi\eps c^3)$. 对 $\pp=p_0\cos\omega t\hat{\vct z}$, 再用 $\langle\cos^2\omega t\rangle=1/2$, 平均角功率与总功率为
\[
\left\langle\frac{\dd P}{\dd\Omega}\right\rangle
=\frac{\mu_0p_0^2\omega^4}{32\pi^2c}\sin^2\theta,\qquad
\langle P\rangle=\frac{p_0^2\omega^4}{12\pi\eps c^3}.
\]
轴向不辐射, 赤道面最大. 对单个非相对论电荷 $\pp=q\rr_p$, 瞬时总功率是拉莫尔公式 $P=q^2a^2/(6\pi\eps c^3)$. 简谐加速度平均平方带来一半, 不能把瞬时和平均系数混用.

\originalsection{近场、感应场和磁偶极}
为处理离源有限距离的题, 保留偶极场的三个距离阶次. 理想时变电偶极在 $r>0$ 的场为
\begin{align*}
\ee={1\over4\pi\eps}\left[
\frac{3\nn(\nn\cdot\pp)-\pp}{r^3}
+\frac{3\nn(\nn\cdot\dot\pp)-\dot\pp}{cr^2}
+\frac{\nn\times(\nn\times\ddot\pp)}{c^2r}\right]_{t-r/c},\\
\bb={\mu_0\over4\pi}\left[\frac{\dot\pp\times\nn}{r^2}
+\frac{\ddot\pp\times\nn}{cr}\right]_{t-r/c}.
\end{align*}
三项依次为近场、感应场、辐射场. 静态极限恢复源外电偶极场. 理想点源的原点接触项不包含在 $r>0$ 公式中.

磁偶极可用真空方程的电磁对偶得到: 以 $\pp=\mm/c$ 的电偶极场作 $\bb_m=\ee_e/c$、$\ee_m=-c\bb_e$. 所以
\begin{align*}
\bb_m={\mu_0\over4\pi}\left[
\frac{3\nn(\nn\cdot\mm)-\mm}{r^3}
+\frac{3\nn(\nn\cdot\dot\mm)-\dot\mm}{cr^2}
+\frac{\nn\times(\nn\times\ddot\mm)}{c^2r}\right]_{t-r/c},\\
\ee_m=-{\mu_0\over4\pi}\left[\frac{\dot\mm\times\nn}{r^2}
+\frac{\ddot\mm\times\nn}{cr}\right]_{t-r/c}.
\end{align*}
磁偶极辐射功率为 $P_m=\mu_0|\ddot\mm|^2/(6\pi c^3)$. 电磁对偶在这里用于真空源外区, 不表示真实材料可以任意交换介电常数和磁导率.

\originalsection{匀速场与相对论辐射}
匀速电荷的加速度项为零. 场相对于电荷\emph{当前}位置差 $\vct R_0$ 可写
\[
\ee=\frac q{4\pi\eps}\frac{(1-\beta^2)\vct R_0}
{[R_0^2-(\vct\beta\times\vct R_0)^2]^{3/2}},\qquad
\bb=c^{-2}\vct v\times\ee.
\]
横向场比静止场强 $\gamma=(1-\beta^2)^{-1/2}$ 倍, 纵向弱 $\gamma^{-2}$ 倍. 这来自推迟几何及洛伦兹变换, 不是匀速运动持续向远方放出辐射.

8.311 将非相对论功率推广为 Li\'enard 公式
\[
P=\frac{q^2\gamma^6}{6\pi\eps c^3}
\left[a^2-\frac{|\vct v\times\vct a|^2}{c^2}\right].
\]
此处作为进阶引用结论, 未声称已从相对论远场角积分完整证明. 平行加速度给 $\gamma^6a^2$, 垂直加速度给 $\gamma^4a^2$, 因而相对论圆周运动的辐射不可用普通拉莫尔公式直接估算.

本章依据 8.07 第16讲、第15讲数学部分及期末题3、4; 源外偶极三阶场用于保留期末题中的有限距离条件. 相对论功率属于向后扩展, 并列明证明范围.
